\documentclass[11pt,a4paper]{article}

\usepackage[margin=2.3cm]{geometry}
\usepackage[T1]{fontenc}
\usepackage{lmodern}
\usepackage{amsmath,amssymb}
\usepackage{booktabs}
\usepackage{array}
\usepackage{tabularx}
\usepackage{xcolor}
\usepackage{enumitem}
\usepackage[hidelinks]{hyperref}

\definecolor{coderef}{HTML}{3775BA}
\definecolor{warn}{HTML}{B64342}
\newcommand{\ctag}[2]{\tag*{\small{\color{coderef}\texttt{#2}}\quad(#1)}}
\newcommand{\code}[1]{\texttt{#1}}
\newcommand{\pmin}{p^{\min}}
\newcommand{\pmax}{p^{\max}}
\newcommand{\pref}{p^{\mathrm{ref}}}
\newcommand{\dPa}{\Delta P^{\mathrm{act}}}
\newcommand{\dPu}{\Delta P^{\mathrm{unb}}}
\newcommand{\fin}{f^{\mathrm{in}}}
\newcommand{\fout}{f^{\mathrm{out}}}
\newcommand{\abar}{\bar a}

\setlength{\parindent}{0pt}
\setlength{\parskip}{4pt}

\title{Pipeline and linepack formulation in AdOpT-NET0\\[2pt]
\large \code{FluidynamicPipeline}: sets, parameters, variables and constraints}
\author{}
\date{10 November 2026}

\begin{document}
\maketitle

\noindent This note writes out the optimization model that the class
\code{FluidynamicPipeline} builds, read off the code at commit \code{9e530a28}.
References in the right margin point to the code:
\code{FP} = \code{fluidynamic\_pipeline.py},
\code{FSP} = \code{fixed\_size\_pipeline.py} (both in
\code{adopt\_net0/components/networks/specificNetworks/}),
\code{NW} = \code{networks/network.py},
\code{FL} = \code{networks/genericNetworks/fluid.py}.
The reasoning behind each modelling choice is in \code{linepack\_formulation.md}
(decisions D1--D11).

Each pipeline type (one json in \code{network\_data/}) is a separate network
object with its own variables. Everything below is written for one such
network; several types at the same node do \emph{not} share variables
(Section~\ref{sec:coupling}).

\tableofcontents

%---------------------------------------------------------------------------
\section{Physics being approximated}

Steady-state, isothermal, single-segment flow in a pipe $e=(i,j)$, hour $t$:
\begin{align}
  f_{e,t}\,\lvert f_{e,t}\rvert &= R_e\,\bigl(p_{i,t}^2 - p_{j,t}^2\bigr)
  && \text{pipeline (Weymouth-type) equation} \label{eq:phys-pipe}\\
  L_{e,t} &= K_e\,\tfrac12\bigl(p_{i,t}+p_{j,t}\bigr)
  && \text{linepack} \label{eq:phys-lp}\\
  L_{e,t} &= L_{e,t-1} + \Delta t\,\bigl(\fin_{e,t}-\fout_{e,t}\bigr),
  \quad f_{e,t} = \tfrac12\bigl(\fin_{e,t}+\fout_{e,t}\bigr)
  && \text{pack balance} \label{eq:phys-bal}
\end{align}
Equations \eqref{eq:phys-lp}--\eqref{eq:phys-bal} are linear.
Only \eqref{eq:phys-pipe} is not; it is replaced by a piecewise-linear (PWL)
relation in Section~\ref{sec:pwl}. Flows are energy flows (MW) and the
linepack is energy (MWh), consistent with the node energy balances.

%---------------------------------------------------------------------------
\section{Sets}

\begin{tabularx}{\linewidth}{@{}l l X@{}}
\toprule
Symbol & Code & Meaning \\
\midrule
$t\in\mathcal T$ & \code{set\_t} & time steps \\
$a=(i,j)\in\mathcal A$ & \code{set\_arcs}, \code{arc\_block[i,j]} &
  directed arcs; a bidirectional pipeline is two arcs $(i,j)$ and $(j,i)$ \\
$\abar=(j,i)$ & & reverse arc of $a$ \\
$e\in\mathcal E$ & \code{set\_arcs\_unique} & pipelines (one per corridor) \\
$n\in\mathcal N^p$ & \code{set\_pressure\_nodes} &
  pressure keys: the node if \code{pressure\_coupled}$=1$, one key per pipeline end if $0$ \\
$z\in\{0,\dots,Z-1\}$ & \code{set\_breakpoints} & breakpoints of the PWL curve \\
\bottomrule
\end{tabularx}

%---------------------------------------------------------------------------
\section{Parameters}

\subsection{Input (json \code{Performance} and topology)}

\begin{tabularx}{\linewidth}{@{}l l X l@{}}
\toprule
Symbol & json key & Meaning & Unit \\
\midrule
$D$, $\varepsilon$ & \code{diameter}, \code{roughness} & pipe diameter, wall roughness & m \\
$\pmin,\ \pmax$ & \code{pressure\_min}, \code{pressure\_max} & pressure band & bar \\
$\pref$ & \code{pressure\_ref} & level at which the PWL curve is built & bar \\
$T$, $Z_c$, $M$ & \code{temperature}, \code{compressibility\_factor}, \code{molar\_mass}
  & gas state & K, --, kg/kmol \\
$ed$ & \code{energy\_density} & energy per mass & MWh/t \\
$Z$ & \code{nr\_breakpoints} & number of breakpoints (default 5) & -- \\
$L_a$ & topology & arc length & km \\
$\gamma_{1..4}$, \dots & economics & capex, opex, losses, compressor coefficients & \\
\bottomrule
\end{tabularx}

\subsection{Derived per arc (\code{fit\_network\_performance}, \code{FSP:78})}

\begin{align}
  R_s &= 8314/M
  \\
  f_D &= \bigl(2\log_{10}(3.7D/\varepsilon)\bigr)^{-2}
  && \text{Nikuradse, fully turbulent}\\
  R_a &= \frac{ed^2\cdot 3.6^2\cdot\frac{\pi^2}{16}\,D^5\cdot 10^{10}}
             {f_D\cdot 1000\cdot Z_c R_s T}\cdot\frac{1}{L_a}
  && [\text{MW}^2/\text{bar}^2]\\
  K_a &= ed\cdot\frac{\pi D^2}{4}\cdot 1000\cdot\frac{100}{Z_c R_s T}\cdot L_a
  && [\text{MWh}/\text{bar}]\\
  B_z &= \frac{z}{Z-1}\,\bigl(\pmax-\pref\bigr), \quad z=0,\dots,Z-1
  && \text{drop breakpoints}\\
  F_{a,z} &= \sqrt{R_a\,B_z\,\bigl(2\pref + B_z\bigr)}
  && \text{flow breakpoints}\\
  \bar F_a &= \max_z F_{a,z}
  && \text{arc capacity (json \code{size\_max} ignored)}
\end{align}
$F_{a,z}$ follows from \eqref{eq:phys-pipe} with the low pressure fixed at
$\pref$ and the high one at $\pref+B_z$ (decision D1). If
\code{linepack\_on}$=0$, then $K_a:=0$ (\code{FP:202}).

\subsection{Derived constants used in the constraints}

\begin{align}
  \dPa &= \max_z B_z = \pmax-\pref, &
  \dPu &= (\pmax-\pmin)-\dPa,
\end{align}
$\Delta t$ = \code{nr\_timesteps\_averaged} (1 without time staging),
$H$ = hours per typical day (0 without typical days),
$\rho$ = \code{rated\_capacity} (1, so sizes and flows are in MW).

%---------------------------------------------------------------------------
\section{Variables}

\begin{tabularx}{\linewidth}{@{}l l l X@{}}
\toprule
Symbol & Code & Domain & Meaning \\
\midrule
\multicolumn{4}{@{}l}{\emph{Per arc $a$}}\\
$y_a$ & \code{var\_installed} & $\{0,1\}$ & arc built \\
$s_a$ & \code{var\_size} & $\ge 0$ & size $=y_a\bar F_a$ \\
$f_{a,t}$ & \code{var\_flow} & $[0,\bar F_a\rho]$ & mean flow in the pipe \\
$\fin_{a,t},\ \fout_{a,t}$ & \code{var\_flow\_in/out} & $[0,\bar F_a\rho]$ & flow entering / leaving \\
$\ell_{a,t}$ & \code{var\_losses} & $[0,\bar F_a\rho]$ & transport losses \\
$\xi_{a,t}$ & \code{var\_direction} & $\{0,1\}$ & direction $a$ active \\
$\delta_{a,t}$ & \code{var\_delta\_pressure} & $[0,\dPa]$ & pressure drop along $a$ \\
$\lambda_{a,t,z}$ & \code{var\_lambda} & $[0,1]$ & PWL weights (SOS2 over $z$) \\
\midrule
\multicolumn{4}{@{}l}{\emph{Per pipeline $e$}}\\
$L_{e,t}$ & \code{var\_linepack} & $[0,K_e\pmax]$ & stored energy (MWh) \\
\midrule
\multicolumn{4}{@{}l}{\emph{Per pressure key $n$}}\\
$p_{n,t}$ & \code{var\_pressure} & $[\pmin,\pmax]$ & pressure (bar) \\
\bottomrule
\end{tabularx}

\medskip
Notation: for an arc $a=(i,j)$, $p_{i,t}$ is the pressure of node $i$
\emph{as seen by arc $a$} (\code{\_get\_pressure}, \code{FP:271}). With
\code{pressure\_coupled}$=1$ all arcs of the network meeting at $i$ see the
same variable. Per pipeline and hour there are 2 binaries ($\xi_a,\xi_{\abar}$)
and 2 SOS2 sets.

%---------------------------------------------------------------------------
\section{Constraints}

All constraints hold for every $t\in\mathcal T$ unless stated otherwise.

\subsection{Investment (per arc)}
\begin{align}
  s_a &= y_a\,\bar F_a \ctag{C1}{FSP:189}\\
  y_a &= 1 \quad\text{if the arc is existing} \ctag{C2}{FSP:187}\\
  \mathrm{capex}_a &= \gamma_1 y_a + \gamma_2 s_a + \gamma_3 L_a y_a + \gamma_4 s_a L_a
    \quad(=0 \text{ if existing}) \ctag{C3}{FSP:210}\\
  y_a &= y_{\abar}, \qquad s_a = s_{\abar} \ctag{C4--C5}{FP:552, NW:858}
\end{align}

\subsection{Flows on an arc}
\begin{align}
  f_{a,t} &= \tfrac12\bigl(\fin_{a,t} + \fout_{a,t}\bigr) \ctag{C6}{FP:374}\\
  \fin_{a,t} &\le \xi_{a,t}\,\bar F_a\rho, \qquad
  \fout_{a,t} \le \xi_{a,t}\,\bar F_a\rho \ctag{C7--C8}{FP:382, FP:392}\\
  f_{a,t} &\le s_a\rho \ctag{C9}{NW:785}\\
  f_{a,t} + f_{\abar,t} &\le s_a\rho \ctag{C10}{NW:871}\\
  \ell_{a,t} &= f_{a,t}\cdot \mathrm{loss}\cdot L_a \ctag{C11}{NW:776}
\end{align}
The GDP disjunction of the parent class (\code{bidirectional\_network\_precise})
is switched off (\code{FP:185}); C7, C8 and C17 enforce the same condition.

\subsection{Pipeline equation: piecewise linear with SOS2 (per arc)}
\label{sec:pwl}
\begin{align}
  \sum_{z} \lambda_{a,t,z} &= \xi_{a,t} \ctag{C12}{FP:706}\\
  \delta_{a,t} &= \sum_{z} B_z\,\lambda_{a,t,z} \ctag{C13}{FP:714}\\
  f_{a,t} &= \sum_{z} F_{a,z}\,\lambda_{a,t,z} \ctag{C14}{FP:724}\\
  \bigl(\lambda_{a,t,0},\dots,\lambda_{a,t,Z-1}\bigr) &\in \mathrm{SOS2} \ctag{C15}{FP:735}
\end{align}
SOS2: at most two \emph{adjacent} weights are non-zero, so $(\delta,f)$ lies on
one segment of the curve. Because the weights sum to $\xi$ (not to 1), an
inactive direction has $\lambda=0$, hence $\delta=0$ and $f=0$. Note that C14
interpolates the \emph{mean} flow $f$.

\subsection{Pressure drop and node pressures (per arc $a=(i,j)$)}
\begin{align}
  \delta_{a,t} &\le (p_{i,t}-p_{j,t}) + \dPa(1-\xi_{a,t}) + \dPu(1-y_a) \ctag{C16a}{FP:675}\\
  \delta_{a,t} &\ge (p_{i,t}-p_{j,t}) - \dPa(1-\xi_{a,t}) - \dPu(1-y_a) \ctag{C16b}{FP:690}
\end{align}
Built and active: $\delta_{a,t}=p_{i,t}-p_{j,t}\ge 0$, gas flows from high to low
pressure. Inactive: $\lvert p_i-p_j\rvert\le\dPa$. Not built: the end pressures
are free over the whole band. The big-M is split in two terms (decision D5)
to keep it as small as possible on built pipelines.

\subsection{Directions and pressures on an idle pipeline (per pipeline $e=(i,j)$)}
\begin{align}
  \xi_{a,t} + \xi_{\abar,t} &\le y_a \ctag{C17}{FP:572}\\
  p_{i,t}-p_{j,t} &\le \dPa\bigl(\xi_{a,t}+\xi_{\abar,t}+1-y_a\bigr) + \dPu(1-y_a) \ctag{C18a}{FP:598}\\
  p_{i,t}-p_{j,t} &\ge -\dPa\bigl(\xi_{a,t}+\xi_{\abar,t}+1-y_a\bigr) - \dPu(1-y_a) \ctag{C18b}{FP:611}
\end{align}
A built pipeline may be idle ($\xi_a=\xi_{\abar}=0$); C18 then forces
$p_{i,t}=p_{j,t}$.

\subsection{Linepack (per pipeline $e=(i,j)$)}
With $\bar p_{e,t}=\tfrac12(p_{i,t}+p_{j,t})$:
\begin{align}
  L_{e,t} &\le K_e\,\pmax\,y_e \ctag{C19}{FP:453}\\
  L_{e,t} &\ge K_e\,\pmin\,y_e \ctag{C20}{FP:466}\\
  L_{e,t} &\le K_e\,\bar p_{e,t} \ctag{C21}{FP:483}\\
  L_{e,t} &\ge K_e\,\bar p_{e,t} - K_e\,\pmax(1-y_e) \ctag{C22}{FP:499}\\
  L_{e,t} &= L_{e,\mathrm{prev}(t)} + \Delta t\sum_{d\in\{a,\abar\}}\bigl(\fin_{d,t}-\fout_{d,t}\bigr) \ctag{C23}{FP:513}
\end{align}
Built: C21--C22 give $L=K\bar p$, i.e.\ \eqref{eq:phys-lp}. Not built: C19 gives
$L=0$. The predecessor is
\[
  \mathrm{prev}(t)=
  \begin{cases}
    t-1, & \text{otherwise},\\
    \lvert\mathcal T\rvert, & t=1 \text{ (no typical days: cyclic over the horizon)},\\
    t+H-1, & (t-1)\bmod H = 0 \text{ (typical days, method 1: cyclic per day)}.
  \end{cases}
\]
No initial linepack is imposed.

\subsection{Coupling to the rest of the model (per node $n$)}
\begin{align}
  \mathrm{inflow}_{n,t} &= \sum_{i\to n}\bigl(\fout_{(i,n),t}-\ell_{(i,n),t}\bigr) \ctag{C24}{FP:753}\\
  \mathrm{outflow}_{n,t} &= \sum_{n\to j}\fin_{(n,j),t} \ctag{C25}{FP:773}\\
  \mathrm{cons}^{\mathrm{send}}_{a,t} &= f_{a,t}\,\bigl(k^{\mathrm{flow}}+k^{\mathrm{flowDist}}L_a\bigr)
  \quad\text{(same at the receiving end)} \ctag{C26}{FL:212}\\
  \mathrm{opex}^{\mathrm{var}} &\propto \textstyle\sum_{a,t} f_{a,t}, \qquad
  \mathrm{opex}^{\mathrm{fix}} \propto \mathrm{capex} \ctag{C27}{NW:812}
\end{align}
The node balance sees the end flows, so gas being packed is not delivered.
Losses, compressor energy and variable opex see the mean flow. The compressor
coefficients $k$ are computed from the fixed json pressures and do not depend
on $p_{n,t}$ (decision D8).

%---------------------------------------------------------------------------
\section{How pressures connect in a network}
\label{sec:coupling}

Three ingredients together act like Kirchhoff's laws:
\begin{enumerate}[leftmargin=*]
  \item \textbf{Current law} (node balance, C24--C25 in the ADOPT energy
    balance): what enters a node equals what leaves it. The reference
    transport model already has this law.
  \item \textbf{Potential}: on every active arc $\delta_a=p_i-p_j$ (C16), with
    $p_i$ one variable per node, shared by every arc of the network at $i$.
  \item \textbf{Pipe law}: $f_a=\mathrm{PWL}(\delta_a)$ (C12--C15).
\end{enumerate}
Around any loop $\sum(p_i-p_j)=0$ identically, so
$\sum_{\text{loop}}\pm\delta_a=0$ and, through the pipe law,
$\sum_{\text{loop}}\pm\mathrm{PWL}^{-1}(f_a)=0$: the \emph{voltage law}. It is
not written explicitly; it follows from point 2. As a consequence, the split of
flow between parallel paths is fixed by the physics instead of being free.

\textbf{Which pressure does an arc see at node $n$?}

\begin{tabularx}{\linewidth}{@{}X l@{}}
\toprule
Arcs meeting at $n$ & Pressures at $n$ \\
\midrule
same pipeline type, \code{pressure\_coupled}$=1$ (default) & one, shared \\
different pipeline types & one per type, not linked \\
same type, \code{pressure\_coupled}$=0$ & one per pipeline end \\
\bottomrule
\end{tabularx}

\smallskip
{\color{warn}Caveat (D9):} \code{medium} and \code{large} share the band
15--30\,bar but get separate pressures at the same node. A fix (R2) is
$p^{\mathrm{medium}}_{n,t}=p^{\mathrm{large}}_{n,t}$ for every shared node.
Between the backbone (30--70\,bar) and the distribution grid the pressures are
independent by design, which acts as a free regulator or compressor.

%---------------------------------------------------------------------------
\section{Tightened variants (\code{examples/ladder/linepack\_formulations.py})}

Applied after the model is built; same integer feasible set, tighter LP
relaxation. $D=\pmax-\pmin$.

\begin{tabularx}{\linewidth}{@{}l l X@{}}
\toprule
Variant & Removes & Adds \\
\midrule
\code{hull} & C21 & $L\le K\bar p - K\pmin(1-y)$ \\
\code{signed} & C16, C18 & $\bigl\lvert(p_i-p_j)-(\delta_a-\delta_{\abar})\bigr\rvert\le D(1-y)$ \\
\code{combined} & both of the above & both of the above \\
\code{combined\_cuts} & as \code{combined} &
  $K\pmin y+\tfrac K2(\delta_a+\delta_{\abar})\le L\le K\pmax y-\tfrac K2(\delta_a+\delta_{\abar})$ \\
\code{magnitude\_cuts} & \code{combined\_cuts}, C12--C15 of both arcs &
  one SOS2 per pipeline: $\sum_z\lambda_{e,t,z}=y$,
  $f_a+f_{\abar}=\sum_z F_z\lambda$, $\delta_a+\delta_{\abar}=\sum_z B_z\lambda$;
  $\xi_a+\xi_{\abar}=y$ with $\xi_{\abar}$ continuous; $\delta_d\le\dPa\xi_d$ \\
\code{magnitude\_ratio\_cuts} & as \code{magnitude\_cuts} &
  $r^{\min}\delta_d\le f_d\le r^{\max}\delta_d$ per direction $d$,
  $r=F_z/B_z$ over $B_z>0$ \\
\bottomrule
\end{tabularx}

%---------------------------------------------------------------------------
\section{Possible extensions (not implemented)}

\paragraph{R1, inequality relaxation.} The curve $F(\Delta p)$ is concave, so
the PWL equals the minimum of its secant lines,
$s_k=\frac{F_{k+1}-F_k}{B_{k+1}-B_k}$, $b_k=F_k-s_kB_k$. Replacing C12--C15 by
\[
  f_{a,t}\le s_k\,\delta_{a,t} + b_k\,\xi_{a,t}\qquad k=0,\dots,Z-2
\]
describes exactly the region \emph{under} the curve without $\lambda$ or SOS2.
It gives up $f\ge F(\Delta p)$, i.e.\ it allows throttling. The direction
binaries $\xi$ must stay: without them both arcs of a pipeline would force
$p_i\ge p_j$ and $p_j\ge p_i$, hence no flow.

\paragraph{Gas velocity limit.} With $\rho=p/(Z_cR_sT)$ and
$\dot m=f/ed$, the velocity is $v=\dot m/(\rho A)$, $A=\pi D^2/4$. A cap
$v\le v^{\max}$ is linear in flow and pressure:
\[
  \fin_{a,t}\le c\,v^{\max}\,p_{i,t},\qquad
  \fout_{a,t}\le c\,v^{\max}\,p_{j,t},\qquad
  c=\frac{ed\cdot A}{Z_cR_sT}\ \text{(with unit conversion)}.
\]
At capacity the current model allows 17--35\,m/s at the inlet and up to
52--81\,m/s if the pipe runs at $\pmin$, so a realistic cap would bind.

%---------------------------------------------------------------------------
\section{What is not modelled}
\begin{itemize}[leftmargin=*]
  \item dependence of the flow on the absolute pressure level (1-D curve at $\pref$, D1);
  \item spatial discretisation, gas inertia, temperature dynamics;
  \item compressor work as a function of the pipeline pressure (D8);
  \item a shared pressure across pipeline types of the same band (D9);
  \item at most one pipeline type per corridor (R3);
  \item an initial linepack (the balance is cyclic).
\end{itemize}

\end{document}
